\documentclass[tikz,border=10pt]{standalone}
\usepackage{ctex}
\usepackage{tikz}
\usetikzlibrary{arrows.meta,positioning,fit,backgrounds}

\begin{document}
\begin{tikzpicture}[
  font=\small,
  >=Latex,
  blk/.style={draw=#1!70!black, fill=#1!10, rounded corners=2.5mm, very thick, minimum width=34mm, minimum height=11mm, align=center},
  sig/.style={->, very thick, draw=black!75},
  box/.style={draw=black!25, rounded corners=3mm, line width=0.8pt}
]

\node[blk=blue] (can) at (0,1.6) {CAN ClockIDS\\时钟倾斜检测};
\node[blk=teal] (ethin) at (0,-1.6) {以太网报文输入\\实时监听线程};
\node[blk=purple] (cpu) at (5.6,0) {中央处理器\\事件融合与告警构建};
\node[blk=orange] (ring) at (11.2,0) {以太网环形存储\\EthRingBuffer};
\node[blk=green!70!black] (alert) at (16.8,0) {告警输出\\含以太网上下文};

\begin{scope}[on background layer]
  \node[box, fit=(can)(ethin)(cpu)(ring)(alert), inner sep=8mm, label={[yshift=2mm]\large\bfseries 车内检测链路中的以太网环形存储位置}] {};
\end{scope}

\draw[sig] (can.east) -- node[above, sloped] {CAN 异常事件} (cpu.west);
\draw[sig] (ethin.east) -- node[below, sloped] {以太网报文流} (cpu.west);
\draw[sig] (cpu.east) -- node[above] {写入/查询请求} (ring.west);
\draw[sig] (ring.east) -- node[above] {时间窗上下文} (alert.west);
\draw[sig, dashed] (cpu.north east) .. controls +(1.6,1.0) and +(-1.2,1.2) .. node[above] {触发检索} (ring.north);

\end{tikzpicture}
\end{document}
